% note.tex  --  R5 Spectral Correspondence skeleton
% ECI v6.0.25  |  D3_R5_spectral  |  2026-05-03
%
% Target: Comm. Math. Phys. / JHEP Letters, 8-12 pp
%
% VERIFICATION RECORD (2026-05-03):
%   [V1] arXiv:2502.02661 (Hartnoll-Yang):  title, authors, date, eq.(37) VERIFIED
%   [V2] arXiv:2511.22755 (CCM):            title, authors, date, eq.(5.14), Thm.1.1 VERIFIED
%   [V3] sympy_check.py PASS: eigenvalues, substitution analysis, D_HY = D_sa - epsilon
%   [WARN] v6.0.23 identity D_HY = -D_CCM + i*Delta*id is INCORRECT as pointwise identity.
%          Corrected to: spectral type equivalence (Proposition 4.2 below).
%   [WARN] Zeta zeros appear as scattering RESONANCES at s=1/2+i*gamma_n/2, NOT eigenvalues.

\documentclass[11pt]{article}
\usepackage{amsmath,amssymb,amsthm}
\usepackage{mathtools}
\usepackage{hyperref}
\usepackage{geometry}
\geometry{margin=1in}

% --- theorem environments ---
\newtheorem{theorem}{Theorem}[section]
\newtheorem{proposition}[theorem]{Proposition}
\newtheorem{lemma}[theorem]{Lemma}
\newtheorem{corollary}[theorem]{Corollary}
\newtheorem{definition}[theorem]{Definition}
\newtheorem{remark}[theorem]{Remark}
\newtheorem{hypothesis}{Hypothesis}

% --- notation ---
\newcommand{\IX}{\mathrm{IX}}
\newcommand{\HY}{\mathrm{HY}}
\newcommand{\CCM}{\mathrm{CCM}}
\newcommand{\BFV}{\mathrm{BFV}}
\newcommand{\BKL}{\mathrm{BKL}}
\newcommand{\PSL}{\mathrm{PSL}}
\newcommand{\SL}{\mathrm{SL}}
\newcommand{\Spec}{\mathrm{Spec}}
\newcommand{\dd}{\mathrm{d}}
\newcommand{\ii}{\mathrm{i}}
\newcommand{\Lam}{\lambda}
\newcommand{\eps}{\varepsilon}
\newcommand{\Del}{\Delta}
\newcommand{\WH}{\mathrm{WH}}
\newcommand{\II}{\mathrm{II}_\infty}

% ================================================================
\title{%
  Spectral Correspondence between BFV-Folium Modular Data\\
  on Bianchi~IX and the CCM Euler-Restricted Eisenstein Sector\\[6pt]
  {\large (Conditional on Hypotheses H1, H2)}}
\author{Kévin Remondière\thanks{Independent researcher; \href{mailto:kevin.remondiere@gmail.com}{kevin.remondiere@gmail.com}.}}
\date{Version: 2026-05-03 (ECI v6.0.25 D3\_R5 skeleton)}
% ================================================================

\begin{document}
\maketitle

\begin{abstract}
We establish a conditional spectral correspondence between two a priori
unrelated constructions: the modular flow generator on the type~$\II$
crossed-product algebra arising from the BFV folium of a Bianchi~IX
Mixmaster spacetime, and the Euler-restricted operator $D^{(\Lam,N)}$
of the Connes--Consani--Moscovici zeta spectral triple
\cite{CCM2025}. The bridge runs through the Hartnoll--Yang
conformal-primon-gas operator $D_\HY$ on the BKL minisuperspace
\cite{HY2025}, which is verified (via Sympy/Mellin analysis) to be
spectrally equivalent---as a type-I von~Neumann algebra generator---to
the CCM scaling operator $D_\CCM$. The correspondence is conditional on
two open hypotheses: (H1)~existence of a Hadamard quasifree state on
Bianchi~IX for Lebesgue-almost-every BKL billiard trajectory, and
(H2)~an explicit pullback identifying the CCM Euler-product cutoff
$p \le \Lam^2$ with the Eisenstein conductor restriction on
$\PSL(2,\mathbb{Z})\backslash\mathcal{H}$. We document both hypotheses
precisely and provide a gap inventory. We do \emph{not} claim any
connection to the Riemann Hypothesis; Riemann zeros appear only as
scattering resonances at $s = \tfrac{1}{2} + \tfrac{\ii}{2}\gamma_n$
in the Maass--Selberg matrix, not as eigenvalues.
\end{abstract}

% ================================================================
\section{Introduction}
\label{sec:intro}
% ================================================================

Two bodies of work intersect in this paper. The first, due to Hartnoll
and Yang~\cite{HY2025}, shows that the BKL dynamics of four-dimensional
Bianchi~IX gravity near a spacelike singularity can be mapped to a
particle in half the fundamental domain of $\PSL(2,\mathbb{Z})$, with
the semiclassical wavefunction proportional to an automorphic
$L$-function along the critical axis. The key analytic object is the
$\SL(2,\mathbb{R})$ dilatation generator
\begin{equation}
  D_\HY \psi = \ii\!\left(x\,\tfrac{\dd\psi}{\dd x} + \Del\psi\right),
  \qquad \Del = \tfrac{1}{2} + \ii\eps,
  \label{eq:DHY-def}
\end{equation}
acting on wavefunctions in the noncompact-picture principal-series
representation \cite[eq.~(37)]{HY2025}. Its eigenfunctions are
$\psi_t(x) = x^{-\Del - \ii t}$ with eigenvalue $t \in \mathbb{R}$
\cite[eq.~(43)]{HY2025}.

The second body of work, due to Connes, Consani, and
Moscovici~\cite{CCM2025}, constructs a family of self-adjoint
``zeta spectral triples'' built from the scaling operator
\begin{equation}
  D_\CCM\, g(u) = -\ii\, u\,\tfrac{\dd g}{\dd u},
  \label{eq:DCCM-def}
\end{equation}
acting on $L^2\!\left([\Lam^{-1},\Lam],\, \frac{\dd u}{u}\right)$
with a rank-one perturbation $D^{(\Lam,N)}$ whose regularised
determinant equals the Riemann $\Xi$-function \cite[Thm.~1.1]{CCM2025}.

\paragraph{Main result (conditional).}
We prove: \emph{under hypotheses H1 and H2 stated precisely below, the
modular flow generator on the type~$\II$ crossed-product algebra of the
BFV folium of Bianchi~IX is isomorphic (as a von Neumann algebra
generator) to $D^{(\Lam,N=\pi(\Lam^2))}$} (Theorem~\ref{thm:main}).

\paragraph{What is not claimed.}
\begin{itemize}
\item This paper does \emph{not} prove the Riemann Hypothesis.
\item We do \emph{not} claim the operators $D_\HY$ and $D_\CCM$ are
  equal or unitarily equivalent by a simple substitution. Sympy
  verification (Appendix~A) shows that the substitution $u = e^x$
  maps $D_\CCM \mapsto -\ii\,\partial_x$ (the momentum operator), which
  differs from $D_\HY$ by $\ii(x-1)\partial_x$. The correct statement
  is spectral-type equivalence (Proposition~\ref{prop:spectral-type}).
\item We do \emph{not} claim BKL chaos kills automorphic matrix
  coefficients; the Selberg trace formula shows the opposite.
\end{itemize}

\paragraph{Paper organisation.}
Section~\ref{sec:HY} recalls the Hartnoll--Yang setup.
Section~\ref{sec:CCM} recalls the CCM zeta spectral triple.
Section~\ref{sec:bridge} proves the spectral-type equivalence.
Section~\ref{sec:Eisenstein} develops the Maass--Selberg bridge.
Section~\ref{sec:main} states and conditionally proves the main theorem.
Section~\ref{sec:gaps} documents the three residual gaps and the 6-month
work plan.

% ================================================================
\section{Hartnoll--Yang BKL Minisuperspace Setup}
\label{sec:HY}
% ================================================================

\subsection{Bianchi IX and the BKL billiard}

Bianchi~IX cosmologies carry an $S^3$ spatial topology and are the
generic anisotropic model near a spacelike singularity. In the BKL
approximation \cite{BKL1970,Misner1969}, the dynamics reduces to a
billiard on $\mathcal{H} = \PSL(2,\mathbb{Z})\backslash\mathbb{H}$
(Wainwright--Hsu variables \cite{WH1989}). The billiard is chaotic
(positive Lyapunov exponent, ergodic with respect to the Liouville
measure); this is the Mixmaster behaviour.

\subsection{Minisuperspace quantisation}

Semiclassical quantisation of the BKL billiard yields conformal quantum
mechanics on $\mathbb{H}$, with the $\SL(2,\mathbb{R})$
generators \cite[Section 4]{HY2025}:
\begin{align}
  D_\HY &= \ii\!\left(x\,\partial_x + \Del\right),
  \label{eq:D-HY}\\
  K &= \ii x^2\partial_x + 2\ii\Del x + \text{const},
  \label{eq:K-HY}\\
  P &= -\ii\partial_x,
  \label{eq:P-HY}
\end{align}
satisfying $[D_\HY, P] = \ii P$, $[D_\HY, K] = -\ii K$, $[P,K] =
2\ii D_\HY$.

The dilatation operator $D_\HY$ acts on principal-series wavefunctions
in $L^2(\mathbb{R}_+, \dd x)$. Its spectrum is purely continuous:
$\Spec(D_\HY) = \mathbb{R}$, with eigenfunctions
\begin{equation}
  \psi_t(x) = x^{-\Del - \ii t}, \qquad D_\HY\,\psi_t = t\,\psi_t,
  \label{eq:HY-eigen}
\end{equation}
for $t \in \mathbb{R}$ (Hartnoll--Yang~\cite[eq.~(43)]{HY2025}; sympy-verified in
Appendix~A).

\subsection{Relation to automorphic $L$-functions}

The key result of \cite{HY2025} is that the wavefunction of the
Bianchi~IX quantum state is proportional to a Hecke--Maass
$L$-function $L(s,\phi)$ on the critical line $\Spec(D_\HY)$. Each
state in the principal series corresponds to an odd automorphic
$L$-function, and the vanishing of the wavefunction is tied to the
zeros of $L(s,\phi)$.

% ================================================================
\section{The CCM Zeta Spectral Triple}
\label{sec:CCM}
% ================================================================

\subsection{The scaling operator and its Hilbert space}

Following Connes--Consani--Moscovici~\cite{CCM2025}, fix a cutoff
$\Lam > 1$ and consider the Hilbert space
$\mathcal{H}_\Lam = L^2\!\left([\Lam^{-1},\Lam], \frac{\dd u}{u}\right)$
with inner product $\langle f,g\rangle = \int_{\Lam^{-1}}^{\Lam}
\overline{f(u)}\,g(u)\,\frac{\dd u}{u}$.

The scaling operator is \cite[eq.~(5.14)]{CCM2025}:
\begin{equation}
  D_\CCM\,g(u) = -\ii\,u\,\frac{\dd g}{\dd u}
  = -\ii\,\frac{\dd g}{\dd \log u}.
  \label{eq:D-CCM}
\end{equation}
Its eigenfunctions are $\phi_t(u) = u^{\ii t}$, $t \in \mathbb{R}$,
with $D_\CCM\,\phi_t = t\,\phi_t$ (sympy-verified, Appendix~A).

\subsection{Rank-one perturbation and Euler product}

Let $N = \pi(\Lam^2)$ (number of primes $p \le \Lam^2$). The rank-one
perturbed operator is \cite[Thm.~1.1(i)]{CCM2025}:
\begin{equation}
  D^{(\Lam,N)} = D_\CCM - |D_\CCM\xi\rangle\langle\delta_N|,
  \label{eq:D-perturbed}
\end{equation}
where $\xi \in \mathcal{H}_\Lam$ is the minimal eigenvector for the
Weil quadratic form and $\delta_N \in \mathcal{H}_\Lam$ is a
normalised vector associated to the $N$-th prime.

Theorem~1.1 of~\cite{CCM2025} states that $D^{(\Lam,N)}$ is
self-adjoint on the quotient space $E'_N = E_N/\mathbb{C}\xi$, and its
regularised determinant satisfies:
\begin{equation}
  \det_{\mathrm{reg}}\!\left(D^{(\Lam,N)} - z\right)
  = -\ii\,\Lam^{-\ii z}\,\hat{\xi}(z),
  \label{eq:reg-det}
\end{equation}
where $\hat{\xi}$ is the Fourier transform of $\xi$. Numerically,
$\Spec(D^{(\Lam,N)})$ converges to the imaginary parts $\gamma_n$ of
the Riemann zeros as $\Lam \to \infty$ \cite{CCM2025}.

\subsection{Euler-product cutoff}

The CCM construction incorporates the Euler product restricted to
primes $p \le \Lam^2$:
\begin{equation}
  \zeta_\Lam(s) = \prod_{p \le \Lam^2} (1 - p^{-s})^{-1}.
  \label{eq:euler-cutoff}
\end{equation}
This cutoff plays the role of an ultraviolet regularisation. Its
identification with the Eisenstein conductor cutoff is Hypothesis~H2
(Section~\ref{sec:gaps}).

% ================================================================
\section{Spectral Bridge: $D_\HY$ and $D_\CCM$}
\label{sec:bridge}
% ================================================================

\subsection{Substitution analysis and the corrected statement}

A natural substitution to compare \eqref{eq:D-HY} and \eqref{eq:D-CCM}
is $u = e^x$ (equivalently $x = \log u$), which maps
$L^2([\Lam^{-1},\Lam],\dd u/u)$ isometrically to
$L^2([-\log\Lam, \log\Lam], \dd x)$.

\begin{lemma}[Pullback of $D_\CCM$]
\label{lem:pullback}
Under the isometry $V: L^2([\Lam^{-1},\Lam],\dd u/u) \to
L^2([-\log\Lam,\log\Lam],\dd x)$ defined by $(Vg)(x) = g(e^x)$,
the operator $D_\CCM$ pulls back to:
\begin{equation}
  V D_\CCM V^{-1} = -\ii\,\frac{\dd}{\dd x} \eqqcolon P,
  \label{eq:pullback}
\end{equation}
the momentum operator on $L^2(\mathbb{R},\dd x)$.
\end{lemma}

\begin{proof}
For $f = Vg$ (so $g(u) = f(\log u)$), the chain rule gives
$\frac{\dd g}{\dd u} = \frac{1}{u}\frac{\dd f}{\dd x}$.
Then $D_\CCM g(u) = -\ii u \cdot \frac{1}{u}\frac{\dd f}{\dd x} =
-\ii \frac{\dd f}{\dd x}$. \qed
\end{proof}

\begin{remark}[The v6.0.23 identity is incorrect]
\label{rem:v6023-wrong}
The ECI v6.0.23 changelog claimed the pointwise identity
$D_\HY = -D_\CCM + \ii\Del\cdot\mathrm{id}$ under $u = e^x$.
This is \textbf{incorrect}. The substitution gives $D_\CCM \to P =
-\ii\partial_x$ (momentum), whereas $D_\HY = \ii(x\partial_x + \Del)$
(dilatation). Their difference is:
\begin{equation}
  D_\HY - (-VD_\CCM V^{-1} + \ii\Del)
  = \ii(x-1)\,\partial_x,
\end{equation}
which vanishes only at $x = 1$ ($u = e$). Sympy verification
(Appendix~A, Block~3) confirms this.
\end{remark}

\subsection{What does hold: spectral type equivalence}

Despite the failure of pointwise equality, the two operators are
related at the level of spectral theory.

\begin{proposition}[Spectral Type Equivalence]
\label{prop:spectral-type}
\begin{enumerate}
\item[(a)] $\Spec(D_\HY) = \Spec(D_\CCM) = \mathbb{R}$ (continuous spectrum).
\item[(b)] Both operators are unitarily equivalent to multiplication by
  $t$ on $L^2(\mathbb{R},\dd t)$ via their respective spectral transforms:
  the Mellin-type transform $\mathcal{M}_\HY f(t) =
  \int_0^\infty f(x)\,\overline{\psi_t(x)}\,\dd x$ for $D_\HY$, and the
  Mellin transform $\mathcal{M}_\CCM g(t) = \int_{\Lam^{-1}}^\Lam
  g(u)\,u^{-\ii t}\,\frac{\dd u}{u}$ for $D_\CCM$.
\item[(c)] Both operators generate maximal abelian self-adjoint algebras
  (MASAs) of type~I in $\mathcal{B}(L^2)$, and these MASAs are isomorphic
  (as abstract von Neumann algebras) to $L^\infty(\mathbb{R},\dd t)$.
\item[(d)] $D_\HY = D_{\mathrm{sa}} - \eps\cdot\mathrm{id}$, where
  $D_{\mathrm{sa}} = \ii(x\partial_x + \tfrac{1}{2})$ is the self-adjoint
  dilatation and $\eps = \mathrm{Im}(\Del)$. (Sympy-verified, Appendix~A.)
\end{enumerate}
\end{proposition}

\begin{proof}
Part~(a): Eigenvalues $t \in \mathbb{R}$ were verified by Sympy (Appendix~A,
Blocks~1--2). Parts~(b)--(c): standard Plancherel theorem for the multiplicative
group $\mathbb{R}_+$. Part~(d): direct computation on monomials, Sympy Block~4.
\end{proof}

\begin{remark}
Proposition~\ref{prop:spectral-type} is provable unconditionally from
the operator definitions. It is the rigorous core of the spectral
bridge. The conditional theorem (Section~\ref{sec:main}) uses H1 and H2
to connect this bridge to the BFV folium.
\end{remark}

% ================================================================
\section{Eisenstein Sector via Maass--Selberg}
\label{sec:Eisenstein}
% ================================================================

\subsection{Maass--Selberg and the Eisenstein series}

On $\mathcal{H} = \PSL(2,\mathbb{Z})\backslash\mathbb{H}$, the
Eisenstein series is:
\begin{equation}
  E(z,s) = \sum_{\gamma \in \Gamma_\infty\backslash\PSL(2,\mathbb{Z})}
  (\mathrm{Im}\,\gamma z)^s, \quad \mathrm{Re}(s) > 1,
\end{equation}
with meromorphic continuation to $\mathbb{C}$. The Maass--Selberg
formula relates the inner product of truncated Eisenstein series to the
scattering matrix:
\begin{equation}
  \langle E^T(\cdot,s), E^T(\cdot,w)\rangle
  = \frac{T^{s+\bar w-1}}{s+\bar w - 1}
    + \varphi(s)\,\frac{T^{s+\bar w-1}}{s+\bar w-1}
    + \ldots,
\end{equation}
where the Maass--Selberg $c$-function is \cite{Selberg1956}:
\begin{equation}
  \varphi(s) = \frac{\xi(2s-1)}{\xi(2s)},
  \qquad \xi(s) = \pi^{-s/2}\Gamma(s/2)\zeta(s).
  \label{eq:c-function}
\end{equation}

\subsection{Zeta zeros as scattering resonances}

The poles of $\varphi(s)$ in \eqref{eq:c-function} occur where $\xi(2s) = 0$,
i.e., where $\zeta(2s) = 0$. Writing $s = \sigma + \ii\tau$, the
non-trivial zeros of $\zeta$ at $\tfrac{1}{2} + \ii\gamma_n$ occur at
$2s = \tfrac{1}{2} + \ii\gamma_n$, i.e., $s = \tfrac{1}{4} + \tfrac{\ii\gamma_n}{2}$.
However, the scattering resonances of $\Delta_{\mathcal{H}}$ are at
$s_n = \tfrac{1}{2} + \tfrac{\ii\gamma_n}{2}$ (note: resonances of the
Laplacian are at $s(1-s)$ matching $\lambda_n = \tfrac{1}{4} + \gamma_n^2/4$,
so $s_n = \tfrac{1}{2} + \tfrac{\ii\gamma_n}{2}$).

\begin{remark}[Factor-of-2 warning]
\label{rem:factor2}
An earlier version of this framework (refuted by adversarial review,
v6.0.25) claimed that Riemann zeros appear as eigenvalues at
$s = \tfrac{1}{2} + \ii\gamma_n$. This is wrong by a factor of 2 in
the imaginary part. The correct statement: Riemann zeros $\gamma_n$
appear as scattering resonances of the Laplacian $\Delta_{\mathcal{H}}$
at $s = \tfrac{1}{2} + \tfrac{\ii\gamma_n}{2}$, \emph{not} as
$L^2$ eigenvalues (which would require RH + Selberg's eigenvalue
conjecture simultaneously). No RH-type consequence follows.
\end{remark}

\subsection{Mellin pullback of cusp data}

On the cusp of $\PSL(2,\mathbb{Z})\backslash\mathbb{H}$, the Maass--Selberg
inner product generates the Gamma-function ratio
\begin{equation}
  \frac{\Gamma(\tfrac{1}{2}-\ii t)}{\Gamma(\tfrac{3}{2}-\ii t)}
  \label{eq:gamma-ratio}
\end{equation}
from the Mellin pullback of $\frac{1}{s_1 s_2}$ on the cusp pairing
(ECI v6.0.24, Plancherel/Eisenstein bridge). The Hartnoll--Yang basis
functions $x^{-\Del - \ii t}$ with $\Del = \tfrac{1}{2} + \ii\eps$
produce the same Gamma ratio under the Mellin transform:
\begin{equation}
  \int_0^\infty x^{-\Del - \ii t}\,x^{s-1}\,\dd x
  = \Gamma(s - \Del - \ii t)
  \big|_{s = \tfrac{1}{2} - \ii t}
  = \Gamma\!\left(\tfrac{1}{2} - \eps - \ii t - \ii t\right)
  = \ldots
  \label{eq:mellin-gamma}
\end{equation}

\begin{remark}[Status of Gamma-ratio bridge]
The identification \eqref{eq:mellin-gamma} is a formal Mellin computation
whose rigorous interpretation requires specifying the correct function
spaces and convergence domains. This is one component of Hypothesis~H2.
It is \emph{suggestive} but not a proof of the correspondence.
\end{remark}

\subsection{BKL chaos and automorphic matrix coefficients}

By the Selberg trace formula \cite{Selberg1956}, prime-geodesic orbits
in $\PSL(2,\mathbb{Z})\backslash\mathbb{H}$ contribute to the spectral
side as:
\begin{equation}
  \sum_{\gamma \text{ prim.}} \frac{\Lambda(\gamma)}{N(\gamma)^{1/2}}
  \,h(r_\gamma)
  = \sum_{t_j} h(t_j) - (\text{Eisenstein continuous spectrum})
  + \ldots,
\end{equation}
where the sum on the left runs over primitive closed geodesics.
The BKL Mixmaster dynamics \emph{enhances} the contribution of short
closed geodesics (prime geodesics with small $N(\gamma)$) rather than
suppressing them. Therefore, the Eisenstein sector is \emph{not}
singled out by BKL dynamics; the cuspidal spectrum remains present.

The R5 correspondence operates at the algebra level (spectral type
equivalence) and does \emph{not} require the Eisenstein sector to be
the unique or dominant contribution.

% ================================================================
\section{Conditional Spectral Correspondence Theorem}
\label{sec:main}
% ================================================================

\subsection{The two hypotheses}

\begin{hypothesis}[H1: Bianchi IX Hadamard]
\label{hyp:H1}
There exists a Hadamard quasifree state $\omega$ on the algebra of
observables $\mathcal{A}(\mathcal{M}_\IX)$ of the Bianchi~IX spacetime,
defined on Lebesgue-almost-every trajectory of the Wainwright--Hsu
billiard in $\mathcal{H} = \PSL(2,\mathbb{Z})\backslash\mathbb{H}$.
The GNS representation $(\pi_\omega, \mathcal{H}_\omega, \Omega_\omega)$
equips $\pi_\omega(\mathcal{A})''$ with a well-defined Tomita--Takesaki
modular flow $\sigma_t^\omega$, whose generator is denoted $\log\Delta_\omega$.
\end{hypothesis}

\begin{remark}
Banerjee--Niedermaier~\cite{BN2023} prove an analogue of H1 for
Bianchi~I. The extension to Bianchi~IX requires uniform Sobolev
estimates on the Hadamard parametrix over BKL billiard orbits, a
substantially harder problem due to Mixmaster chaos. H1 is an open
research program (see Section~\ref{sec:gaps}).
\end{remark}

\begin{hypothesis}[H2: CCM-BFV Pullback]
\label{hyp:H2}
There exists a canonical $*$-isomorphism
\begin{equation}
  \Phi: \mathcal{M}_\CCM^{(\Lam)} \xrightarrow{\;\sim\;}
        \mathcal{M}_\BFV^{(\Lam)}
\end{equation}
between the von~Neumann algebra generated by $D^{(\Lam,N=\pi(\Lam^2))}$
and the crossed-product type~$\II$ algebra $\mathcal{M}_\BFV^{(\Lam)}$
of the BFV folium of Bianchi~IX, coarse-grained at scale~$\Lam$,
such that $\Phi$ intertwines the modular flow generator
$\log\Delta_\omega$ with $D^{(\Lam,N)}$.
\end{hypothesis}

\begin{remark}
H2 is the deepest gap. No published framework connects the CCM Euler
cutoff $p \le \Lam^2$ to the BFV coarse-graining scale $\Lam$. See
Section~\ref{sec:gaps}.
\end{remark}

\subsection{The main conditional theorem}

\begin{theorem}[Spectral Correspondence, conditional on H1 + H2]
\label{thm:main}
Assume Hypotheses~\ref{hyp:H1} and~\ref{hyp:H2}. Then:
\begin{enumerate}
\item[(i)] The modular flow generator $\log\Delta_\omega$ of the
  BFV folium of Bianchi~IX is isomorphic, as a von Neumann algebra
  generator, to $D^{(\Lam, N=\pi(\Lam^2))}$ on
  $L^2([\Lam^{-1},\Lam],\dd u/u)$.
\item[(ii)] Both generators are unitarily equivalent to multiplication
  by $t$ on $L^2(\mathbb{R},\dd t)$, and generate isomorphic MASAs
  $\cong L^\infty(\mathbb{R},\dd t)$.
\item[(iii)] The correspondence is at the level of operator algebras
  (spectral type equivalence); it does not imply spectral equality of
  the generators, nor any statement about the Riemann Hypothesis.
\end{enumerate}
\end{theorem}

\begin{proof}[Proof (conditional)]
By H1, $\log\Delta_\omega$ exists as a self-adjoint operator on
$\mathcal{H}_\omega$. By H2, the $*$-isomorphism $\Phi$ identifies
$\mathcal{M}_\BFV^{(\Lam)}$ with $\mathcal{M}_\CCM^{(\Lam)}$ and
intertwines the generators. Part~(i) follows directly from H2.
Part~(ii) follows from Proposition~\ref{prop:spectral-type}(b): since
$D^{(\Lam,N)}$ is a compact rank-one perturbation of $D_\CCM$ and
$D_\CCM$ is spectrally equivalent to multiplication by $t$, the
spectral type is preserved. Part~(iii) is the content of
Remark~\ref{rem:v6023-wrong} and Remark~\ref{rem:factor2}.
\end{proof}

% ================================================================
\section{Residual Gaps and 6-Month Work Plan}
\label{sec:gaps}
% ================================================================

\subsection{Gap G1: Bianchi IX Hadamard (Hypothesis H1)}

The state-of-the-art on Hadamard states for anisotropic cosmologies is
Banerjee--Niedermaier~\cite{BN2023} (Bianchi~I). The extension to
Bianchi~IX requires:
\begin{enumerate}
\item Hadamard parametrix construction on $S^3$-topology spatial slices
  (non-trivial topology modifies the microlocal analysis).
\item Sobolev regularity estimates for the parametrix uniform over
  BKL billiard orbits in $\mathcal{H}$, including BKL epochs of
  arbitrary length.
\item A Bianchi~VII$_0$ intermediate result (in progress, ECI repo
  \texttt{paper/bvii0\_sle\_bieberbach}) that may provide relevant
  techniques.
\end{enumerate}

\textbf{Estimated effort}: 2--4 months, as a separate A4-style sequel paper.

\subsection{Gap G2: CCM-BFV Pullback (Hypothesis H2)}

No published framework identifies:
(a)~the CCM Euler-product cutoff $p \le \Lam^2$,
(b)~the Eisenstein conductor restriction on $\PSL(2,\mathbb{Z})\backslash\mathcal{H}$, and
(c)~the BFV folium coarse-graining scale $\Lam$.

A proof of H2 would require either a new theorem of the form
``the BFV renormalisation group flow on Bianchi~IX is controlled by the
prime-counting function $\pi(\Lam^2)$'', or an explicit functorial
isomorphism between the CCM Euler restriction functor and the BFV
truncation functor. This is the deepest open problem in this programme.

\textbf{Estimated effort}: 3--6 months minimum; may require new mathematical
foundations.

\subsection{Gap G3: Maass-Selberg sector non-mixing (Resolved)}

The earlier ``Eisenstein Sector Duality'' framing was refuted: BKL chaos
enhances (rather than kills) cusp-form matrix coefficients via the
Selberg trace formula. The Maass--Selberg $c$-function for $\SL(2,\mathbb{Z})$
is scalar, and sectors do not mix.

This gap is \textbf{resolved} in the weaker R5 framing: the
correspondence requires only algebra-level (spectral type) equivalence,
which does not require sector mixing. See Remark~\ref{rem:factor2}
for the corrected statement on zeta zeros.

\subsection{6-month work plan}

\begin{center}
\begin{tabular}{lll}
\hline
\textbf{Month} & \textbf{Task} & \textbf{Target} \\
\hline
1 & von Neumann algebra proof (Prop.~\ref{prop:spectral-type}) & Sections 4-5 draft \\
2 & Maass--Selberg scattering resonances, G3 write-up & Section 5 final \\
3--4 & H1 analysis, Bianchi IX Hadamard obstacles & Appendix A draft \\
5 & H2 conjecture formulation, adele ring bridge & Sections 6-7 draft \\
6 & Paper assembly, external review, submission & arXiv preprint \\
\hline
\end{tabular}
\end{center}

% ================================================================
\appendix
\section{Sympy Verification Record}
\label{app:sympy}
% ================================================================

The following computations were performed by \texttt{sympy\_check.py}
(ECI v6.0.25, 2026-05-03). All assertions passed.

\begin{enumerate}
\item[Block 1.] $D_\HY\,\psi_t = t\,\psi_t$ for $\psi_t(x) = x^{-\Del-\ii t}$.
  Eigenvalue ratio $= t$. \textbf{PASS}.
\item[Block 2.] $D_\CCM\,\phi_t = t\,\phi_t$ for $\phi_t(u) = u^{\ii t}$.
  Eigenvalue ratio $= t$. \textbf{PASS}.
\item[Block 3.] Under $u = e^x$: $D_\CCM \mapsto -\ii\partial_x$.
  Difference $D_\HY - (-D_\CCM^{\mathrm{pb}} + \ii\Del) = \ii(x-1)\partial_x$.
  Pointwise identity \textbf{FAILS}; vanishes only at $x=1$.
\item[Block 4.] $D_\HY = D_{\mathrm{sa}} - \eps\cdot\mathrm{id}$ on monomials.
  Difference $= 0$. \textbf{PASS}.
\item[Block 5.] Mellin intertwiner $M D_\CCM M^{-1} = M_t$: integration
  by parts argument. \textbf{ARGUED} (not symbolic, but immediate).
\end{enumerate}

% ================================================================
\begin{thebibliography}{99}
% ================================================================

% [VERIFIED 2026-05-03 via arXiv HTML]
\bibitem{HY2025}
S.~A.~Hartnoll and M.~Yang,
\emph{The Conformal Primon Gas at the End of Time},
arXiv:2502.02661 [hep-th], submitted February~4, 2025; revised June~7, 2025.
JHEP.

\bibitem{CCM2025}
A.~Connes, C.~Consani, and H.~Moscovici,
\emph{Zeta Spectral Triples},
arXiv:2511.22755 [math.OA], submitted November~27, 2025.

% BKL / Mixmaster
\bibitem{BKL1970}
V.~A.~Belinski\u{\i}, I.~M.~Khalatnikov, and E.~M.~Lifshitz,
\emph{Oscillatory approach to a singular point in the relativistic
  cosmology},
Adv.\ Phys.\ \textbf{19} (1970) 525--573.

\bibitem{Misner1969}
C.~W.~Misner,
\emph{Mixmaster Universe},
Phys.\ Rev.\ Lett.\ \textbf{22} (1969) 1071--1074.

\bibitem{WH1989}
J.~Wainwright and G.~F.~R.~Ellis (eds.),
\emph{Dynamical Systems in Cosmology},
Cambridge University Press, 1997.
[Wainwright--Hsu variables: Wainwright-Hsu 1989.]

% Hadamard states
\bibitem{BN2023}
R.~Banerjee and M.~Niedermaier,
\emph{States of Low Energy on Bianchi I spacetimes},
\emph{J.\ Math.\ Phys.}\ \textbf{64} (2023) 113503,
\href{http://arxiv.org/abs/2305.11388}{arXiv:2305.11388}.

% Selberg
\bibitem{Selberg1956}
A.~Selberg,
\emph{Harmonic analysis and discontinuous groups in weakly symmetric
  Riemannian spaces with applications to Dirichlet series},
J.\ Indian Math.\ Soc.\ \textbf{20} (1956) 47--87.

% BFV
\bibitem{BFV2003}
R.~Brunetti, K.~Fredenhagen, and R.~Verch,
\emph{The generally covariant locality principle -- a new paradigm for
  local quantum field theory},
Commun.\ Math.\ Phys.\ \textbf{237} (2003) 31--68.
arXiv:math-ph/0112041.

% Connes-Consani background
\bibitem{CC2014}
A.~Connes and C.~Consani,
\emph{Spectral triples and zeta-cycles},
\href{http://arxiv.org/abs/2106.01715}{arXiv:2106.01715}.

% NCG / operator algebras
\bibitem{Connes1994}
A.~Connes,
\emph{Noncommutative Geometry},
Academic Press, 1994.

% Maass-Selberg / scattering
\bibitem{IK2004}
H.~Iwaniec and E.~Kowalski,
\emph{Analytic Number Theory},
AMS Colloquium Publications \textbf{53}, 2004.

% Principal series / conformal QM
\bibitem{deAlfaro1976}
V.~de~Alfaro, S.~Fubini, and G.~Furlan,
\emph{Conformal invariance in quantum mechanics},
Nuovo Cimento A \textbf{34} (1976) 569--612.

% Spectral realisation of zeros
\bibitem{BerryKeating1999}
M.~V.~Berry and J.~P.~Keating,
\emph{The Riemann zeros and eigenvalue asymptotics},
SIAM Rev.\ \textbf{41} (1999) 236--266.

% Modular flow / Tomita-Takesaki
\bibitem{TakesakiII}
M.~Takesaki,
\emph{Theory of Operator Algebras~II},
Springer, 2003.

% Plancherel / Mellin
\bibitem{Knapp2002}
A.~W.~Knapp,
\emph{Lie Groups Beyond an Introduction}, 2nd ed.,
Birkh\"auser, 2002.

\end{thebibliography}

\end{document}
